\documentclass[aps,pra,preprint,nofootinbib]{revtex4-2}

\usepackage{amsmath,amssymb,amsthm,bm,mathtools}
\usepackage{microtype}
\usepackage{hyperref}
\usepackage{booktabs}

\newtheorem{theorem}{Theorem}
\newtheorem{proposition}{Proposition}
\newtheorem{lemma}{Lemma}
\newtheorem{corollary}{Corollary}
\newtheorem{definition}{Definition}
\newtheorem{remark}{Remark}

\newcommand{\Tr}{\operatorname{Tr}}
\newcommand{\supp}{\operatorname{supp}}
\newcommand{\id}{\mathbb I}
\newcommand{\Rlin}{R_{\mathrm{lin}}}
\newcommand{\Aact}{\mathcal A}

\begin{document}

\title{Supplemental Material for ``Two spectral-resource regimes for autonomous temporal information''}
\author{Anonymous}
\affiliation{Anonymous}
\maketitle

This Supplemental Material gives self-contained proof cores for the resource statements used in the main text.  The order follows the physical hierarchy: finite-radius survival, rank-changing boundary synthesis, the arbitrary-support mixed bridge, the multi-gap shared-Hessian law, and finally the exact qutrit common-record accessibility benchmark.  Standard ingredients---numerical-radius inequalities, PSD-cone second-order geometry, shorted operators, and multiparameter Fisher bounds---are used as mathematical infrastructure rather than novelty claims \cite{BhatiaJain2023,Shapiro1997,BonnansCominettiShapiro1999,AndersonTrapp1975,BraunsteinCaves1994,GillMassar2000}.

\section{Finite-radius survival}
% M2 proof module: WP02 / WP03 / WP06 finite-radius resource laws.

\subsection{Finite-radius arbitrary-POVM survival laws}
\label{supp:finite-radius-full}

We use the convention
\begin{equation}
D_c=(A+A^\dagger)/2,
\qquad
D_s=(A-A^\dagger)/(2i).
\end{equation}
Let $P=\supp\rho_0$ and suppose the affine disk
\begin{equation}
\rho_0+xD_c+yD_s\succeq0
\qquad(x^2+y^2\le R^2)
\end{equation}
is physical for some $R>0$. Then positivity in both signs implies $A=PAP$.

Define on $P$
\begin{equation}
B=\rho_0^{-1/2}A\rho_0^{-1/2}.
\end{equation}
Writing $\beta=x-iy$, the affine family is congruent to
\begin{equation}
\id+\frac12(\beta B+\beta^*B^\dagger).
\end{equation}
For a unit vector $|\psi\rangle$, minimizing the expectation value over the phase of $\beta$ gives
\begin{equation}
1-|\beta|\,|\langle\psi|B|\psi\rangle|.
\end{equation}
Therefore the exact affine radius is
\begin{equation}
R_{\rm lin}=\frac1{w(B)}.
\label{eq:proof-Rlin}
\end{equation}
The standard numerical-radius inequality $\|B\|\le2w(B)$ gives
\begin{equation}
B B^\dagger\preceq\frac4{R_{\rm lin}^2}\id_P.
\end{equation}
By congruence,
\begin{equation}
A\rho_0^+A^\dagger
=\rho_0^{1/2}BB^\dagger\rho_0^{1/2}
\preceq\frac4{R_{\rm lin}^2}\rho_0.
\label{eq:proof-operator-majorization}
\end{equation}
Similarly,
\begin{equation}
A^\dagger\rho_0^+A
\preceq\frac4{R_{\rm lin}^2}\rho_0.
\end{equation}

For any POVM $\{M_y\}$ define
\begin{equation}
p_y=\Tr(\rho_0M_y),
\qquad
z_y=\Tr(AM_y).
\end{equation}
Weighted Hilbert--Schmidt Cauchy--Schwarz gives
\begin{equation}
\frac{|z_y|^2}{p_y}
\le\Tr(M_yA\rho_0^+A^\dagger).
\end{equation}
After summing the outcomes,
\begin{equation}
\Tr F_1
\le\Tr(A\rho_0^+A^\dagger).
\label{eq:proof-score-bound}
\end{equation}

Now suppose the range of $A$ lies in a projector $P_U$: $P_UA=A$. The positive operator on the left of Eq.~\eqref{eq:proof-operator-majorization} is then supported in $P_U$, hence
\begin{align}
\Tr F_1
&\le\Tr(A\rho_0^+A^\dagger)\\
&=\Tr[P_U(A\rho_0^+A^\dagger)]\\
&\le\frac4{R_{\rm lin}^2}\Tr(P_U\rho_0).
\end{align}
Thus
\begin{equation}
\frac{R_{\rm lin}^2}{4}\Tr F_1
\le\Tr(P_U\rho_0).
\label{eq:proof-upper-tail}
\end{equation}
The same argument applied to $A^\dagger$ yields a domain/lower-endpoint bound.

\paragraph{Finite-copy extension.}
For $N$ independently encoded copies,
\begin{equation}
A_N=\sum_{j=1}^N
\rho_0^{\otimes(j-1)}\otimes A\otimes
\rho_0^{\otimes(N-j)}.
\end{equation}
Because $\Tr A=0$, all cross-copy terms in
\begin{equation}
\Tr(A_N\rho_N^+A_N^\dagger)
\end{equation}
vanish, while every diagonal copy term equals the one-copy weighted norm. Therefore
\begin{equation}
\Tr(A_N\rho_N^+A_N^\dagger)
=N\Tr(A\rho_0^+A^\dagger).
\end{equation}
The weighted score inequality is valid for an arbitrary POVM on the full $N$-copy Hilbert space, so
\begin{equation}
\frac{R_{\rm lin}^2}{4}\frac{\Tr F_N}{N}
\le\Tr(P_U\rho_0).
\label{eq:proof-upper-tail-N}
\end{equation}
No product-measurement assumption is used.

\paragraph{Exact Bohr mode.}
If
\begin{equation}
[H,A_\nu]=\hbar\nu A_\nu,
\end{equation}
then the range and domain of $A_\nu$ lie in upper and lower endpoint sectors separated by $\hbar\nu$. For a stationary baseline, applying Eq.~\eqref{eq:proof-upper-tail-N} to $A_\nu$ and $A_\nu^\dagger$ gives
\begin{equation}
\frac{R_{\rm lin}^2}{4}\frac{\Tr F_N^{(\nu)}}{N}
\le\min\{U_\nu,D_\nu\}.
\end{equation}
For a semibounded Hamiltonian the upper endpoint lies inside the tail
\begin{equation}
T(\nu)=\Pr_{\rho_0}[H-E_*\ge\hbar\nu],
\end{equation}
so
\begin{equation}
\frac{R_{\rm lin}^2}{4}\frac{\Tr F_N^{(\nu)}}{N}
\le T(\nu).
\end{equation}
Markov's inequality then gives the mean-energy corollary.

\subsection{Autonomous dual survival}
\label{supp:dual-survival-full}

For the exact clock--signal exchange
\begin{equation}
[H_S,A_\nu]=+\hbar\nu A_\nu,
\qquad
[H_C,A_\nu]=-\hbar\nu A_\nu,
\end{equation}
the signal viewpoint sees $A_\nu$ as a positive $+\nu$ tangent, whereas the clock viewpoint sees $A_\nu^\dagger$ as a positive $+\nu$ tangent. The two Hermitian tangent quadratures span the same physical plane under $A_\nu\leftrightarrow A_\nu^\dagger$; only the sine coordinate reverses. Hence the affine radius and Fisher trace are unchanged.

Applying Eq.~\eqref{eq:proof-upper-tail-N} to the signal range projector and to the clock range projector of $A_\nu^\dagger$ yields
\begin{equation}
\frac{R_{\rm lin}^2}{4}\frac{\Tr F_N^{(\nu)}}{N}
\le T_S(\nu),
\end{equation}
and
\begin{equation}
\frac{R_{\rm lin}^2}{4}\frac{\Tr F_N^{(\nu)}}{N}
\le T_C(\nu).
\end{equation}
Therefore
\begin{equation}
\boxed{
\frac{R_{\rm lin}^2}{4}\frac{\Tr F_N^{(\nu)}}{N}
\le\min\{T_C(\nu),T_S(\nu)\}.
}
\end{equation}
The proof requires only global physicality of the baseline and exact local gap structure of the tangent; separate local stationarity is unnecessary because Eq.~\eqref{eq:proof-upper-tail-N} did not assume $[\rho_0,H_X]=0$.


\section{Rank-changing boundary synthesis and autonomous action}
% M2 proof module: WP07 / WP09 / audited WP18.

\subsection{Rank-changing boundary curvature and autonomous action}
\label{supp:boundary-action-full}

Let $\rho(x,y)$ be a $C^2$ physical family with baseline $\rho_0$, support projector $P=\supp\rho_0$, and kernel projector $Q=\id-P$. Let
\begin{equation}
\partial_x\rho(0)=D_c=\frac{A+A^\dagger}{2},
\qquad
\partial_y\rho(0)=D_s=\frac{A-A^\dagger}{2i}.
\end{equation}

\subsubsection{One-sided boundary synthesis}

Assume a baseline-empty endpoint projector $P_U$ with
\begin{equation}
A=P_UAP,
\qquad
P_UP=0.
\end{equation}
Define
\begin{equation}
J=\Tr(A\rho_0^+A^\dagger),
\qquad
T_U(x,y)=\Tr[P_U\rho(x,y)].
\end{equation}
For any POVM, the right-supported weighted score inequality gives
\begin{equation}
\Tr F_1\le J.
\label{eq:boundary-score-one}
\end{equation}
The independent-copy identity gives
\begin{equation}
\frac{\Tr F_N}{N}\le J
\label{eq:boundary-score-N}
\end{equation}
for arbitrary collective measurements.

To obtain the curvature cost, apply the second-order PSD-cone condition separately to the $x$ and $y$ curves. Since
\begin{equation}
QD_cP=\frac A2,
\qquad
QD_sP=\frac A{2i},
\end{equation}
one has
\begin{align}
Q\partial_x^2\rho(0)Q
&\succeq
2\left(\frac A2\right)\rho_0^+
\left(\frac A2\right)^\dagger
=\frac12A\rho_0^+A^\dagger,\\
Q\partial_y^2\rho(0)Q
&\succeq
2\left(\frac A{2i}\right)\rho_0^+
\left(\frac A{2i}\right)^\dagger
=\frac12A\rho_0^+A^\dagger.
\end{align}
Adding gives
\begin{equation}
C_\Delta
:=Q(\partial_x^2+\partial_y^2)\rho(0)Q
\succeq A\rho_0^+A^\dagger.
\label{eq:one-side-Cdelta}
\end{equation}
Taking the trace after projection onto $P_U$ yields
\begin{equation}
\Delta T_U(0)
\ge J.
\end{equation}
Thus
\begin{equation}
\boxed{
\frac{\Tr F_N}{N}\le J\le\Delta T_U(0).
}
\label{eq:one-side-final}
\end{equation}

\subsubsection{Bilateral boundary Minkowski law}

Now assume a pure-boundary tangent
\begin{equation}
PAP=0,
\qquad
QAQ=0.
\end{equation}
Define
\begin{equation}
X=QAP,
\qquad
Y=QA^\dagger P.
\end{equation}
Then
\begin{equation}
A=X+Y^\dagger.
\end{equation}
For a POVM define
\begin{equation}
u_y=\frac{\Tr(XM_y)}{\sqrt{p_y}},
\qquad
v_y=\frac{\Tr(YM_y)^*}{\sqrt{p_y}}.
\end{equation}
The Fisher trace is the squared norm of the sum of score vectors,
\begin{equation}
\Tr F_1=\|u+v\|_2^2.
\end{equation}
Weighted Cauchy--Schwarz gives
\begin{equation}
\|u\|_2^2\le J_+,
\qquad
\|v\|_2^2\le J_-,
\end{equation}
where
\begin{equation}
J_+=\Tr(X\rho_0^+X^\dagger),
\qquad
J_-=\Tr(Y\rho_0^+Y^\dagger).
\end{equation}
Therefore Minkowski gives
\begin{equation}
\boxed{
\sqrt{\Tr F_1}\le\sqrt{J_+}+\sqrt{J_-}.
}
\end{equation}
For $N$ copies the two weighted norms scale as $NJ_+$ and $NJ_-$, hence
\begin{equation}
\boxed{
\sqrt{\frac{\Tr F_N}{N}}
\le\sqrt{J_+}+\sqrt{J_-}.
}
\label{eq:bilateral-N}
\end{equation}

For the second-order resource, the support-to-kernel blocks of the two real derivatives are
\begin{equation}
K_c=\frac{X+Y}{2},
\qquad
K_s=\frac{X-Y}{2i}.
\end{equation}
The two PSD-cone inequalities give
\begin{equation}
Q\partial_x^2\rho Q\succeq2K_c\rho_0^+K_c^\dagger,
\end{equation}
\begin{equation}
Q\partial_y^2\rho Q\succeq2K_s\rho_0^+K_s^\dagger.
\end{equation}
Adding them cancels the $X/Y$ cross terms exactly:
\begin{equation}
C_\Delta\succeq
X\rho_0^+X^\dagger+Y\rho_0^+Y^\dagger
=Z_++Z_-.
\label{eq:bilateral-shared-curvature}
\end{equation}

If $X$ and $Y$ enter orthogonal baseline-empty endpoint sectors $P_+$ and $P_-$, define
\begin{equation}
T_\pm(x,y)=\Tr(P_\pm\rho(x,y)).
\end{equation}
Projecting Eq.~\eqref{eq:bilateral-shared-curvature} onto the two sectors gives
\begin{equation}
J_+\le\Delta T_+(0),
\qquad
J_-\le\Delta T_-(0).
\end{equation}
Combining with Eq.~\eqref{eq:bilateral-N},
\begin{equation}
\frac{\Tr F_N}{N}
\le
\left[\sqrt{\Delta T_+(0)}+\sqrt{\Delta T_-(0)}\right]^2.
\label{eq:bilateral-curvature-final}
\end{equation}
The cross term in the square is generally physical and cannot be replaced by an additive curvature bound.

\subsubsection{Autonomous two-sided action}

Assume now the exact exchange relations
\begin{equation}
[H_S,A_\nu]=+\hbar\nu A_\nu,
\qquad
[H_C,A_\nu]=-\hbar\nu A_\nu.
\end{equation}
In a clean bilateral boundary geometry, $X$ creates a signal-upper and clock-lower endpoint, whereas $Y$ creates a signal-lower and clock-upper endpoint. Let the corresponding local population Laplacians be
\begin{equation}
\Delta T_{S,+},\quad\Delta T_{S,-},\quad
\Delta T_{C,+},\quad\Delta T_{C,-}.
\end{equation}
Define
\begin{equation}
\Aact_S^{(2)}
=\frac{\hbar\nu}{4}
(\Delta T_{S,+}+\Delta T_{S,-}),
\end{equation}
\begin{equation}
\Aact_C^{(2)}
=\frac{\hbar\nu}{4}
(\Delta T_{C,+}+\Delta T_{C,-}).
\end{equation}
From Eq.~\eqref{eq:bilateral-curvature-final} and
\begin{equation}
(\sqrt a+\sqrt b)^2\le2(a+b),
\end{equation}
each subsystem obeys
\begin{equation}
\Aact_X^{(2)}
\ge\frac{\hbar\nu}{8}\frac{\Tr F_N}{N},
\qquad X\in\{C,S\}.
\end{equation}
Therefore
\begin{equation}
\boxed{
\Aact_C^{(2)}+\Aact_S^{(2)}
\ge\frac{\hbar\nu}{4}\frac{\Tr F_N}{N}.
}
\label{eq:autonomous-bilateral-action-proof}
\end{equation}

If only one orientation is present, Eq.~\eqref{eq:one-side-final} applies on each local side separately. Each subsystem then obeys
\begin{equation}
\Aact_X^{(2)}
\ge\frac{\hbar\nu}{4}\frac{\Tr F_N}{N},
\end{equation}
so
\begin{equation}
\boxed{
\Aact_C^{(2)}+\Aact_S^{(2)}
\ge\frac{\hbar\nu}{2}\frac{\Tr F_N}{N}.
}
\label{eq:autonomous-one-sided-action-proof}
\end{equation}

\subsubsection{Sharp fixed-shell constructions}

For bilateral synthesis use the total-excitation-2 shell
\begin{equation}
|L\rangle=|2_C,0_S\rangle,
\quad
|M\rangle=|1_C,1_S\rangle,
\quad
|U\rangle=|0_C,2_S\rangle,
\end{equation}
with
\begin{equation}
A_\nu=c(|U\rangle\langle M|+|M\rangle\langle L|),
\qquad
\rho_0=|M\rangle\langle M|.
\end{equation}
The exact family
\begin{align}
|\psi(x,y)\rangle={}&
\sqrt{1-\frac{c^2}{2}(x^2+y^2)}|M\rangle\\
&+\frac c2(x+i y)|L\rangle
+\frac c2(x-i y)|U\rangle
\end{align}
is normalized on the open disk $x^2+y^2<2/c^2$ and lies entirely inside one total-energy eigenspace. Every local endpoint Laplacian is $c^2$, while the three-outcome Fourier measurement gives
\begin{equation}
\Tr F_1=4c^2.
\end{equation}
Thus Eq.~\eqref{eq:autonomous-bilateral-action-proof} is saturated exactly.

For one-sided synthesis use
\begin{equation}
A_\nu=2c|U\rangle\langle D|,
\qquad
\rho_0=|D\rangle\langle D|,
\end{equation}
inside a fixed-total-excitation-1 shell and the convention-consistent family
\begin{equation}
|\psi(x,y)\rangle
=\sqrt{1-c^2(x^2+y^2)}|D\rangle+c(x-i y)|U\rangle.
\end{equation}
The local endpoint Laplacians are $4c^2$, a fixed equatorial POVM gives $\Tr F_1=4c^2$, and Eq.~\eqref{eq:autonomous-one-sided-action-proof} is saturated exactly.


\section{Arbitrary coherent support}
% M2 proof module: audited WP11/WP13/WP19 arbitrary-support autonomous bridge.

\subsection{Arbitrary coherent-support mixed survival--synthesis bridge}
\label{supp:mixed-full-proof}

Let $A_\nu$ satisfy the exact exchange relations
\begin{equation}
[H_S,A_\nu]=+\hbar\nu A_\nu,
\qquad
[H_C,A_\nu]=-\hbar\nu A_\nu,
\end{equation}
and let $P=\supp\rho_0$, $Q=\id-P$. A two-sided differentiable physical family implies
\begin{equation}
QA_\nu Q=0.
\end{equation}
Decompose
\begin{equation}
B=PA_\nu P,
\qquad
K_+=QA_\nu P,
\qquad
K_-=QA_\nu^\dagger P.
\label{eq:mixed-proof-decomp}
\end{equation}
Then
\begin{equation}
A_\nu=B+K_++K_-^\dagger.
\end{equation}
Define
\begin{align}
J_B^+&=\Tr(B\rho_0^+B^\dagger),&
J_B^-&=\Tr(B^\dagger\rho_0^+B),\\
J_+&=\Tr(K_+\rho_0^+K_+^\dagger),&
J_-&=\Tr(K_-\rho_0^+K_-^\dagger).
\end{align}

\subsubsection{Measurement-side bound}

Set
\begin{equation}
X=(B+K_+)P,
\qquad
Y=K_-.
\end{equation}
The outputs of $B$ and $K_+$ lie in orthogonal $P$ and $Q$ sectors, hence
\begin{equation}
\Tr(X\rho_0^+X^\dagger)=J_B^++J_+.
\end{equation}
For a POVM, the complex score is the sum of the right-supported $X$ score and the conjugate $Y$ score. Weighted Hilbert--Schmidt Cauchy--Schwarz and Minkowski therefore give
\begin{equation}
\sqrt{\Tr F_1}
\le
\sqrt{J_B^++J_+}+\sqrt{J_-}.
\end{equation}
Applying the same argument to $A_\nu^\dagger$ gives
\begin{equation}
\sqrt{\Tr F_1}
\le
\sqrt{J_B^-+J_-}+\sqrt{J_+}.
\end{equation}
The independent-copy weighted norms scale by $N$, so for every finite $N$ and arbitrary collective POVM,
\begin{equation}
\boxed{
\sqrt{\frac{\Tr F_N}{N}}
\le
\min\!\left\{
\sqrt{J_B^++J_+}+\sqrt{J_-},
\sqrt{J_B^-+J_-}+\sqrt{J_+}
\right\}.
}
\label{eq:mixed-proof-measurement}
\end{equation}

\subsubsection{Two-sided shorted survival ceilings}

Because $P$ need not commute with either local Hamiltonian, $B=PAP$ is generally not itself an exact Bohr-gap operator. The energy resource of $B$ therefore depends on the angle between its information-bearing support and the physical endpoint subspaces.

Let
\begin{equation}
Z_B^+=B\rho_0^+B^\dagger,
\qquad
Z_B^-=B^\dagger\rho_0^+B,
\end{equation}
and
\begin{equation}
R_B^+=\supp Z_B^+,
\qquad
R_B^-=\supp Z_B^-.
\end{equation}
Let $\Pi_{S,U},\Pi_{S,D},\Pi_{C,U},\Pi_{C,D}$ be the minimal local endpoint projectors containing the participating support-preserving transition content. Define compressed endpoint operators
\begin{align}
S_{S,U}&=P\Pi_{S,U}P,&S_{S,D}&=P\Pi_{S,D}P,\\
S_{C,U}&=P\Pi_{C,U}P,&S_{C,D}&=P\Pi_{C,D}P.
\end{align}
The four shorting constants are
\begin{align}
\lambda_{S,U}&=\sup\{\lambda:S_{S,U}\succeq\lambda R_B^+\},\\
\lambda_{C,D}&=\sup\{\lambda:S_{C,D}\succeq\lambda R_B^+\},\\
\lambda_{S,D}&=\sup\{\lambda:S_{S,D}\succeq\lambda R_B^-\},\\
\lambda_{C,U}&=\sup\{\lambda:S_{C,U}\succeq\lambda R_B^-\}.
\end{align}
Let $R_B$ be the physical affine radius of the support-preserving sub-tangent $B$ and $T_{X,E}=\Tr(\Pi_{X,E}\rho_0)$. The finite-radius operator majorization gives
\begin{align}
J_B^+&\le\frac{4T_{S,U}}{R_B^2\lambda_{S,U}},&
J_B^+&\le\frac{4T_{C,D}}{R_B^2\lambda_{C,D}},\\
J_B^-&\le\frac{4T_{S,D}}{R_B^2\lambda_{S,D}},&
J_B^-&\le\frac{4T_{C,U}}{R_B^2\lambda_{C,U}}.
\end{align}
Thus define
\begin{align}
a_+&=
\min\!\left\{
\frac{4T_{S,U}}{R_B^2\lambda_{S,U}},
\frac{4T_{C,D}}{R_B^2\lambda_{C,D}}
\right\},\\
a_-&=
\min\!\left\{
\frac{4T_{S,D}}{R_B^2\lambda_{S,D}},
\frac{4T_{C,U}}{R_B^2\lambda_{C,U}}
\right\},
\end{align}
so
\begin{equation}
J_B^+\le a_+,
\qquad
J_B^-\le a_-.
\label{eq:mixed-proof-apm}
\end{equation}
If $B=0$, set $a_+=a_-=0$.

\subsubsection{Canonical endpoint-incidence action}

Define the canonical joint endpoint-role projectors of the full exact exchange,
\begin{equation}
\Pi_{\rm out}=\supp(A_\nu A_\nu^\dagger),
\qquad
\Pi_{\rm in}=\supp(A_\nu^\dagger A_\nu).
\end{equation}
The exchange commutators imply
\begin{equation}
[H_X,A_\nu A_\nu^\dagger]=0,
\qquad
[H_X,A_\nu^\dagger A_\nu]=0
\end{equation}
for $X=C,S$, hence the two support projectors commute with both local Hamiltonians. They encode the output and input endpoint roles of the full ladder.

One joint endpoint incidence corresponds to one absolute signal gap and one absolute clock gap. Define
\begin{equation}
G_{\rm ex}
=2\hbar\nu\,Q(\Pi_{\rm out}+\Pi_{\rm in})Q.
\label{eq:mixed-proof-Gex}
\end{equation}
If a joint energy state is both a domain and range endpoint, it is counted twice because it plays two distinct roles.

Let
\begin{equation}
C_\Delta=Q(\partial_x^2+\partial_y^2)\rho(0)Q.
\end{equation}
Second-order positivity gives
\begin{equation}
C_\Delta\succeq Z_++Z_-,
\end{equation}
where
\begin{equation}
Z_+=K_+\rho_0^+K_+^\dagger,
\qquad
Z_-=K_-\rho_0^+K_-^\dagger.
\end{equation}
Define
\begin{equation}
\Aact_{\rm ex}^{(2)}
=\frac14\Tr(G_{\rm ex}C_\Delta).
\end{equation}
Let $R_\pm=\supp Z_\pm$ and
\begin{equation}
g_\pm=
\lambda_{\min}
(R_\pm G_{\rm ex}R_\pm|_{R_\pm})
\end{equation}
for nonzero components. Then
\begin{align}
4\Aact_{\rm ex}^{(2)}
&=\Tr(G_{\rm ex}C_\Delta)\\
&\ge\Tr(G_{\rm ex}Z_+)+\Tr(G_{\rm ex}Z_-)\\
&\ge g_+J_+ +g_-J_-.
\label{eq:mixed-proof-action-budget}
\end{align}
The same physical curvature is charged once.

\subsubsection{Exact scalar reduction}

For $a,e\ge0$ and $p,q>0$, maximize
\begin{equation}
\left[\sqrt{a+j_+}+\sqrt{j_-}\right]^2
\end{equation}
subject to
\begin{equation}
pj_+ +qj_-\le e,
\qquad j_+,j_-\ge0.
\end{equation}
Because the objective is increasing, the budget is saturated at the optimum. Set
\begin{equation}
j_-=\frac{e-pj_+}{q},
\qquad0\le j_+\le e/p.
\end{equation}
The derivative of the square-root sum vanishes when
\begin{equation}
\frac1{\sqrt{a+j_+}}
=\frac{p/q}{\sqrt{(e-pj_+)/q}},
\end{equation}
which yields an interior solution exactly when
\begin{equation}
e\ge\frac{ap^2}{q}.
\end{equation}
Evaluating the boundary and interior branches gives
\begin{equation}
\Psi_a(e;p,q)=
\begin{cases}
(\sqrt a+\sqrt{e/q})^2,&e\le ap^2/q,\\[1mm]
(e+pa)(p^{-1}+q^{-1}),&e\ge ap^2/q.
\end{cases}
\label{eq:mixed-proof-Psi}
\end{equation}

Use Eq.~\eqref{eq:mixed-proof-apm} and Eq.~\eqref{eq:mixed-proof-action-budget} in the first orientation of Eq.~\eqref{eq:mixed-proof-measurement} to obtain
\begin{equation}
\frac{\Tr F_N}{N}
\le
\Psi_{a_+}(4\Aact_{\rm ex}^{(2)};g_+,g_-).
\end{equation}
The conjugate orientation gives
\begin{equation}
\frac{\Tr F_N}{N}
\le
\Psi_{a_-}(4\Aact_{\rm ex}^{(2)};g_-,g_+).
\end{equation}
Therefore
\begin{equation}
\boxed{
\frac{\Tr F_N}{N}
\le
\min\!\left\{
\Psi_{a_+}(4\Aact_{\rm ex}^{(2)};g_+,g_-),
\Psi_{a_-}(4\Aact_{\rm ex}^{(2)};g_-,g_+)
\right\}.
}
\label{eq:mixed-proof-final}
\end{equation}

When $K_+=K_-=0$, the action term vanishes and Eq.~\eqref{eq:mixed-proof-final} reduces to two-sided finite-radius survival. When $B=0$ and the endpoint geometry is clean, $g_+=g_-=2\hbar\nu$ and the bilateral harmonic price is $\hbar\nu$, recovering the sharp autonomous boundary coefficient $\hbar\nu/4$.

\subsubsection{Shared-kernel qutrit check}

For the fixed-total-excitation-2 qutrit benchmark,
\begin{equation}
A=|M\rangle\langle L|-\sqrt2|U\rangle\langle M|,
\end{equation}
so
\begin{equation}
\Pi_{\rm out}=|M\rangle\langle M|+|U\rangle\langle U|,
\end{equation}
\begin{equation}
\Pi_{\rm in}=|L\rangle\langle L|+|M\rangle\langle M|.
\end{equation}
Thus
\begin{equation}
G_{\rm ex}/(\hbar\nu)=\operatorname{diag}(2,4,2).
\end{equation}
The exact weighted norms are
\begin{equation}
J_B^+=J_B^-=\frac54,
\qquad
J_+=\frac74,
\qquad
J_-=3.
\end{equation}
Both synthesized ranges are the same rank-one kernel direction, giving
\begin{equation}
g_+=g_-=\frac{13}{4}\hbar\nu.
\end{equation}
For the minimal curvature,
\begin{equation}
4\Aact_{\rm ex}^{(2)}=\frac{247}{16}\hbar\nu.
\end{equation}
Using $a_+=a_-=5/4$ and units $\hbar\nu=1$, the second branch of Eq.~\eqref{eq:mixed-proof-Psi} gives
\begin{equation}
\Psi_{5/4}(247/16;13/4,13/4)=12.
\end{equation}
Hence the canonical autonomous action reduction introduces no additional resource-layer looseness in this benchmark.


\section{Multi-gap shared-Hessian action}
% M2 proof module: audited WP20 multi-gap shared-Hessian action sum.

\subsection{Multi-gap shared-Hessian spectral-action sum}
\label{supp:multigap-full-proof}

Let $\rho(\bm x,\bm y)$ be one common $C^2$ physical family with baseline $\rho_0$, support projector $P$, and kernel projector $Q$. For each positive gap $\nu_k$, let the complex tangent be $A_k$ and assume a pure-boundary decomposition
\begin{equation}
P A_k P=0,
\qquad
Q A_k Q=0.
\end{equation}
Define
\begin{equation}
X_k=Q A_k P,
\qquad
Y_k=Q A_k^\dagger P,
\end{equation}
so
\begin{equation}
A_k=X_k+Y_k^\dagger.
\end{equation}
The cosine and sine support-to-kernel blocks are
\begin{equation}
K_{k,c}=\frac{X_k+Y_k}{2},
\qquad
K_{k,s}=\frac{X_k-Y_k}{2i}.
\end{equation}

For the $x_k$ coordinate, second-order positivity gives
\begin{equation}
Q\partial_{x_k}^2\rho(0)Q
\succeq2K_{k,c}\rho_0^+K_{k,c}^\dagger.
\end{equation}
For the $y_k$ coordinate,
\begin{equation}
Q\partial_{y_k}^2\rho(0)Q
\succeq2K_{k,s}\rho_0^+K_{k,s}^\dagger.
\end{equation}
Adding the two inequalities cancels the mixed $X_k/Y_k$ terms exactly:
\begin{align}
C_{\Delta,k}
&:=Q(\partial_{x_k}^2+\partial_{y_k}^2)\rho(0)Q\\
&\succeq
X_k\rho_0^+X_k^\dagger
+Y_k\rho_0^+Y_k^\dagger\\
&=:Z_{k,+}+Z_{k,-}.
\label{eq:mode-shared-curvature-proof}
\end{align}
Now sum the operator inequalities over all mode coordinates before introducing any scalar resource:
\begin{equation}
\boxed{
C_\Sigma
:=Q\sum_k(\partial_{x_k}^2+\partial_{y_k}^2)\rho(0)Q
\succeq
\sum_k(Z_{k,+}+Z_{k,-}).
}
\label{eq:common-hessian-proof}
\end{equation}
No mixed second derivatives are needed because $C_\Sigma$ is the trace of the parameter Hessian in the chosen orthonormal coordinate system.

Let $G\succeq0$ be one positive kernel cost operator and define
\begin{equation}
\Aact_{G,\Sigma}^{(2)}
=\frac14\Tr(GC_\Sigma).
\end{equation}
For each nonzero orientation let
\begin{equation}
g_{k,+}
=\lambda_{\min}
(R_{k,+}GR_{k,+}|_{R_{k,+}}),
\end{equation}
\begin{equation}
g_{k,-}
=\lambda_{\min}
(R_{k,-}GR_{k,-}|_{R_{k,-}}),
\end{equation}
where $R_{k,\pm}=\supp Z_{k,\pm}$. From Eq.~\eqref{eq:common-hessian-proof},
\begin{equation}
4\Aact_{G,\Sigma}^{(2)}
\ge
\sum_k(g_{k,+}J_{k,+}+g_{k,-}J_{k,-}),
\label{eq:total-cost-proof}
\end{equation}
with $J_{k,\pm}=\Tr Z_{k,\pm}$.

For one fixed POVM on the $N$ copies, the bilateral boundary Minkowski law applies to each Fisher block $F_{N,k}$ of that same record:
\begin{equation}
\sqrt{\frac{\Tr F_{N,k}}{N}}
\le
\sqrt{J_{k,+}}+\sqrt{J_{k,-}}.
\end{equation}
For positive bilateral prices, define the harmonic price
\begin{equation}
\gamma_k
=\left(g_{k,+}^{-1}+g_{k,-}^{-1}\right)^{-1}.
\end{equation}
Weighted Cauchy--Schwarz gives
\begin{equation}
\gamma_k
(\sqrt{J_{k,+}}+\sqrt{J_{k,-}})^2
\le
g_{k,+}J_{k,+}+g_{k,-}J_{k,-}.
\end{equation}
For a one-sided mode, set $\gamma_k$ equal to its one nonzero orientation cost. Hence
\begin{equation}
\gamma_k\frac{\Tr F_{N,k}}{N}
\le
g_{k,+}J_{k,+}+g_{k,-}J_{k,-}.
\end{equation}
Summing and using Eq.~\eqref{eq:total-cost-proof},
\begin{equation}
\boxed{
\sum_k\gamma_k\frac{\Tr F_{N,k}}{N}
\le4\Aact_{G,\Sigma}^{(2)}.
}
\label{eq:multigap-proof-final}
\end{equation}
This theorem does not require one measurement to be individually optimal for all modes. It applies to the Fisher blocks of any one common record.

\subsubsection{Clean autonomous frequency weighting}

Suppose each mode is an exact clock--signal exchange at gap $\nu_k$ and the clean upper/lower endpoint ranges are mode separated. Choose $G$ to assign the positive combined clock-plus-signal incidence price
\begin{equation}
2\hbar\nu_k
\end{equation}
to either orientation of mode $k$. Then
\begin{equation}
g_{k,+}=g_{k,-}=2\hbar\nu_k,
\end{equation}
so
\begin{equation}
\gamma_k=\hbar\nu_k.
\end{equation}
Equation~\eqref{eq:multigap-proof-final} becomes
\begin{equation}
\boxed{
\Aact_{G,\Sigma}^{(2)}
\ge
\sum_k\frac{\hbar\nu_k}{4}
\frac{\Tr F_{N,k}}{N}.
}
\label{eq:clean-multigap-proof-final}
\end{equation}

\subsubsection{One common sharp Fourier measurement}

Take the fixed-total-excitation shell
\begin{equation}
|n\rangle
\equiv|m-n\rangle_C|m+n\rangle_S,
\qquad n=-m,\ldots,m,
\end{equation}
and baseline $|0\rangle$. For $k=1,\ldots,m$ choose
\begin{equation}
A_k=c_k(|k\rangle\langle0|+|0\rangle\langle-k|),
\qquad
\nu_k=k\omega_0.
\end{equation}
The exact common family is
\begin{align}
|\psi(\bm x,\bm y)\rangle={}&
\sqrt{1-\frac12\sum_{k=1}^m c_k^2(x_k^2+y_k^2)}|0\rangle\\
&+\sum_{k=1}^m\frac{c_k}{2}(x_k+i y_k)|-k\rangle
+\sum_{k=1}^m\frac{c_k}{2}(x_k-i y_k)|k\rangle.
\end{align}
It is physical on the open ellipsoid
\begin{equation}
\frac12\sum_kc_k^2(x_k^2+y_k^2)<1
\end{equation}
and lies entirely in one total-energy eigenspace. Its kernel Hessian is
\begin{equation}
C_\Sigma
=\sum_{k=1}^m c_k^2
\left(|-k\rangle\langle-k|+|k\rangle\langle k|\right).
\end{equation}

Let $d=2m+1$ and measure the discrete Fourier basis
\begin{equation}
|v_j\rangle
=\frac1{\sqrt d}\sum_{n=-m}^{m}e^{in\phi_j}|n\rangle,
\qquad
\phi_j=\frac{2\pi j}{d}.
\end{equation}
At the baseline $p_j=1/d$. For mode $k$,
\begin{equation}
z_{j,k}
=\langle v_j|A_k|v_j\rangle
=\frac{2c_k}{d}e^{-ik\phi_j}.
\end{equation}
Therefore the cosine and sine score components are discrete harmonics. Since $1\le k,l\le m$ and $d=2m+1$,
\begin{align}
\sum_{j=0}^{d-1}\cos(k\phi_j)\cos(l\phi_j)
&=\frac d2\delta_{kl},\\
\sum_{j=0}^{d-1}\sin(k\phi_j)\sin(l\phi_j)
&=\frac d2\delta_{kl},\\
\sum_{j=0}^{d-1}\cos(k\phi_j)\sin(l\phi_j)
&=0.
\end{align}
Thus the full common-record Fisher matrix is
\begin{equation}
F_{x_kx_l}=2c_k^2\delta_{kl},
\qquad
F_{y_ky_l}=2c_k^2\delta_{kl},
\qquad
F_{x_ky_l}=0.
\end{equation}
In particular,
\begin{equation}
\Tr F_{1,k}=4c_k^2
\end{equation}
for every $k$ simultaneously.

Choose
\begin{equation}
G=2\hbar\omega_0
\sum_{n\ne0}|n|\,|n\rangle\langle n|.
\end{equation}
Then
\begin{equation}
\Aact_{G,\Sigma}^{(2)}
=\sum_{k=1}^m\hbar\nu_k c_k^2
\end{equation}
while
\begin{equation}
\sum_{k=1}^m
\frac{\hbar\nu_k}{4}\Tr F_{1,k}
=\sum_{k=1}^m\hbar\nu_k c_k^2.
\end{equation}
Hence one common measurement simultaneously saturates every mode and the complete weighted sum.

For overlapping ranges in a general model, Eq.~\eqref{eq:multigap-proof-final} remains valid for any supplied $G\succeq0$, but the existence of a simple frequency-diagonal physical interpretation becomes an operator-design question. No uniqueness claim is made outside the clean mode-separated geometry.


\section{Exact common-record accessibility benchmark}
% Proof module for the exact common-record Fisher supremum in the shared-kernel qutrit.
% Intended for later inclusion in the Supplemental Material.

\subsection{Exact common-record Fisher supremum for the shared-kernel qutrit}
\label{supp:wp15-exact}

Use the fixed-total-excitation qutrit benchmark defined in the main supplement. Choose an orthonormal basis $\{|e_1\rangle,|e_2\rangle\}$ for the support of $\rho_0$ such that, in the ordered basis $\{|e_1\rangle,|e_2\rangle,|q\rangle\}$,
\begin{equation}
\rho_0=\operatorname{diag}(1/2,1/2,0)
\end{equation}
and
\begin{equation}
A=
\begin{pmatrix}
0 & \sqrt2/2 & \sqrt{30}/6\\
\sqrt2/4 & 0 & \sqrt6/3\\
\sqrt{30}/12 & -\sqrt6/3 & 0
\end{pmatrix}.
\label{eq:wp15-A}
\end{equation}
For a POVM effect $M_y$, let
\begin{equation}
p_y=\Tr(\rho_0M_y),\qquad z_y=\Tr(AM_y).
\end{equation}
The two-quadrature classical Fisher trace is
\begin{equation}
\Tr F_1=\sum_{y:p_y>0}\frac{|z_y|^2}{p_y}.
\end{equation}

\paragraph{Rank-one refinement.}
If $M=\sum_jM_j$ is decomposed into positive effects and $p_j=\Tr(\rho_0M_j)$, $z_j=\Tr(AM_j)$, then quadratic-over-linear convexity gives
\begin{equation}
\frac{|\sum_jz_j|^2}{\sum_jp_j}
\le\sum_j\frac{|z_j|^2}{p_j}.
\end{equation}
Therefore refining effects into rank-one components cannot decrease the Fisher trace. It is sufficient to establish a pointwise bound for $M=w|\phi\rangle\langle\phi|$.

\paragraph{Dual quadratic witness.}
Define
\begin{equation}
Y=
\begin{pmatrix}
9/16 & 0 & 0\\
0 & 9/16 & 3\sqrt{15}/8\\
0 & 3\sqrt{15}/8 & 23/4
\end{pmatrix}.
\label{eq:wp15-Y}
\end{equation}
We claim
\begin{equation}
|\langle\phi|A|\phi\rangle|^2
\le
\langle\phi|\rho_0|\phi\rangle
\langle\phi|Y|\phi\rangle
\label{eq:wp15-witness}
\end{equation}
for every vector $|\phi\rangle$.

For $\lambda>0$ and phase $\theta$, define
\begin{equation}
\mathcal M(\lambda,\theta)
=
\lambda\rho_0+\lambda^{-1}Y
-\left(e^{i\theta}A+e^{-i\theta}A^\dagger\right).
\label{eq:wp15-LMI}
\end{equation}
Set $x=\lambda^2$ and $t=\cos^2\theta$. Direct symbolic reduction gives the first leading principal minor
\begin{equation}
\frac{8x+9}{16\lambda}>0,
\end{equation}
the second leading principal minor
\begin{equation}
\frac{
 t(8x-9)^2+(1-t)(64x^2+112x+81)
}{256x},
\label{eq:wp15-minor2}
\end{equation}
and the determinant
\begin{equation}
\frac{
 t(8x-9)^2+(1-t)(40x+81)
}{128\lambda^3}.
\label{eq:wp15-det}
\end{equation}
For $x>0$ and $0\le t\le1$, both numerators are manifestly nonnegative. They are strictly positive except at the isolated point $t=1$, $x=9/8$; positive semidefiniteness there follows by continuity. Hence
\begin{equation}
\mathcal M(\lambda,\theta)\succeq0
\qquad\text{for all }\lambda>0,\ \theta\in\mathbb R.
\end{equation}
For an arbitrary $|\phi\rangle$, choose $\theta$ to align the phase of $\langle\phi|A|\phi\rangle$. Then
\begin{equation}
2|\langle\phi|A|\phi\rangle|
\le
\lambda\langle\phi|\rho_0|\phi\rangle
+\lambda^{-1}\langle\phi|Y|\phi\rangle.
\end{equation}
Minimizing over $\lambda$ proves Eq.~\eqref{eq:wp15-witness}.

For a rank-one POVM, Eq.~\eqref{eq:wp15-witness} gives
\begin{equation}
\Tr F_1
\le
\sum_y w_y\langle\phi_y|Y|\phi_y\rangle
=\Tr Y
=\frac{55}{8}.
\label{eq:wp15-upper}
\end{equation}
Rank-one refinement extends the upper bound to every POVM.

\paragraph{Sharp limiting projective sequence.}
Write Eq.~\eqref{eq:wp15-A} in block form
\begin{equation}
A=\begin{pmatrix}B&b\\a^\dagger&0\end{pmatrix},
\end{equation}
where
\begin{equation}
B=
\begin{pmatrix}
0&\sqrt2/2\\
\sqrt2/4&0
\end{pmatrix},
\quad
 a=\left(\frac{\sqrt{30}}{12},-\frac{\sqrt6}{3}\right),
\quad
 b=\begin{pmatrix}\sqrt{30}/6\\\sqrt6/3\end{pmatrix}.
\end{equation}
The numerical radius of $B$ is
\begin{equation}
w(B)=\frac{3\sqrt2}{8}.
\end{equation}
The two orthogonal support vectors
\begin{equation}
|s_\pm\rangle=\frac{|e_1\rangle\pm|e_2\rangle}{\sqrt2}
\end{equation}
therefore contribute a total
\begin{equation}
\frac98
\end{equation}
to the Fisher trace.

For a vector approaching the kernel,
\begin{equation}
|\phi_\epsilon\rangle
=\cos\epsilon|q\rangle+\sin\epsilon|t\rangle,
\end{equation}
the baseline probability is $\sin^2\epsilon/2$, while its Fisher contribution tends to
\begin{equation}
2|\langle q|A|t\rangle+\langle t|A|q\rangle|^2.
\end{equation}
Choosing
\begin{equation}
|t_*\rangle=i\frac{(a-b)^T}{\|a-b\|}
\end{equation}
gives
\begin{equation}
\|a-b\|^2=\frac{23}{8},
\end{equation}
so the near-kernel contribution tends to
\begin{equation}
\frac{23}{4}.
\end{equation}
Let $U_\epsilon$ be a unitary rotation by angle $\epsilon$ in $\operatorname{span}\{|q\rangle,|t_*\rangle\}$ and the identity on its orthogonal complement. The projective measurement
\begin{equation}
\left\{
U_\epsilon|s_+\rangle,
U_\epsilon|s_-\rangle,
U_\epsilon|q\rangle
\right\}
\end{equation}
has ordinary positive baseline probabilities for every sufficiently small $\epsilon>0$, and
\begin{equation}
\lim_{\epsilon\to0^+}\Tr F_1[U_\epsilon]
=\frac98+\frac{23}{4}
=\frac{55}{8}.
\label{eq:wp15-lower}
\end{equation}
Combining Eqs.~\eqref{eq:wp15-upper} and \eqref{eq:wp15-lower},
\begin{equation}
\boxed{
\sup_{\mathrm{POVM}}\Tr F_1=\frac{55}{8}.
}
\end{equation}
The result is a supremum; exact attainment by a regular POVM is not required.

For this benchmark the three distinct layers are therefore
\begin{equation}
12\quad>\quad\frac{43}{4}\quad>\quad\frac{55}{8},
\end{equation}
corresponding respectively to the physical shared-curvature resource ceiling, the SLD-QFI trace, and the exact accessible one-copy common-record Fisher supremum.


\section{Interpretive separation of resource, QFI, and accessibility}
The inequalities proved above are physical necessary-resource bounds on an encoded state family.  They are not SLD-attainability statements.  Conversely, a QFI matrix is a quantum-statistical tangent bound and need not equal the physical resource ceiling.  The shared-kernel qutrit makes this separation explicit:
\begin{equation}
12>\frac{43}{4}>\frac{55}{8},
\end{equation}
where the three quantities are respectively the physical resource ceiling, the SLD-QFI trace, and the exact one-copy common-record Fisher supremum.  This distinction is important at rank-changing boundaries, where zero-probability directions and incompatible score quadratures can separate the three layers substantially \cite{Safranek2017,LiEtAl2016}.

\bibliographystyle{apsrev4-2}
\bibliography{references}

\end{document}
